\section{The Question Left by Being Given}

Once a human being grasps that existence was never owed to him, the ordinary arithmetic of obligation breaks down. Between equals, debt has a beginning and a clear end: someone lends a sum or renders a service, and the borrower restores the same measure and walks away clear. Creation admits no such symmetry. Before God called the world out of nothing, there was no creditor waiting in the dark to be paid, no prior contract, and no capital standing on the creature's side of the ledger. Everything a human being might offer back to God---breath, speech, intellect, choice, time, and the very earth beneath his feet---belongs already to the giver. St.\ Paul halts his long exposition of divine mercy in Romans at precisely that precipice: ``Or who hath first given to him, and recompense shall be made him? For of him, and by him, and in him, are all things: to him be glory for ever'' (Rom 11:35--36).\endnote{Scripture quotations in English follow the Challoner revision of the Douay-Rheims Bible; Latin Scripture follows the Clementine Vulgate. On the relationship between Douay-Rheims/Vulgate and Hebrew Psalm numbering, see note 2 below.}

Yet the impossibility of antecedent claim does not leave the creature silent or unoccupied. Immediately after that doxology, St.\ Paul pivots from wonder to concrete bodily offering: ``I beseech you therefore, brethren, by the mercy of God, that you present your bodies a living sacrifice, holy, pleasing unto God, your reasonable service'' (Rom 12:1). The Latin phrase \emph{rationabile obsequium}---``reasonable service''---names a worship shaped by mind and truth rather than mechanical compulsion. It raises the practical question left by the awe of unowed existence. If being itself is sheer gift and can never be repaid in kind, in what sense is worship owed in justice at all? How does a due return take embodied, liturgical, moral, and merciful form even while human affection remains cold or divided? And why can no fulfilled act, however faithful, ever close the creature's account?

The Psalter puts the difficulty into an unforgettable question: \emph{Quid retribuam Domino pro omnibus quae retribuit mihi?}---``What shall I render to the Lord, for all the things that he hath rendered to me?'' (Ps 115:12).\endnote{In the Challoner Douay-Rheims Bible and the liturgical Vulgate, Hebrew Psalm 116 is divided into two psalms: Psalm 114 (Heb.\ 116:1--9) and Psalm 115 (Heb.\ 116:10--19). Challoner continues the verse numbers from 10 to 19, giving Ps 115:12--13 and 115:16--17; the Clementine Vulgate and the 1899 American Douay-Rheims renumber Psalm 115 from 1 to 10, making these verses 115:3--4 and 115:7--8 (corresponding throughout to Hebrew Ps 116:12--13 and 116:16--17).} St.\ Augustine, preaching on this verse in his \emph{Exposition on Psalm 115}, stops dead on the verb itself. The psalmist does not speak of benefits merely given or done, but uses \emph{retribuit}---benefits \emph{returned}. Augustine presses the word until its irony breaks open:

\begin{quote}
He saith not, for all the benefits that He hath done unto me but ``for all the benefits that He hath returned unto me.'' What deeds then on the man's part had preceded, that all the benefits of God were not said to be given, but returned? What had preceded, on the man's part, save sins? God therefore repayeth good for evil, whilst unto Him men repay evil for good; for such was the return of those who said, ``This is the heir: come, let us kill him.''\endnote{St.\ Augustine, \emph{Enarrationes in Psalmos}, Ps 115, \S 9 (numbered Ps CXVI in the Nicene and Post-Nicene Fathers English translation, Series I, vol.\ 8, where the editor notes ``Lat.\ CXV'').}
\end{quote}

The psalmist's dilemma is therefore sharper than mere metaphysical smallness. Even an unfallen creature would stand before God with empty hands, since every capacity to give is itself received (1 Cor 4:7). A fallen creature stands worse than empty-handed: what preceded God's mercy was rebellion. God repays evil with gift, and the creature is left asking what conceivable return could match a generosity that precedes every desert and outruns every sin.

\subsection{The Reluctant Son and the Worth of a Due Act}

Many serious believers quietly suspect that an act performed out of duty---before warm affection or spontaneous delight has caught up with it---is hollow, insincere, or even hypocritical. If I go to Mass when my heart is dull, pray the words of the psalms while distracted, or keep a precept while feeling inward reluctance, am I merely performing an empty shell that God despises?

In St.\ Matthew's Gospel, Christ refuses to collapse honesty into spontaneity. He puts a brief parable to the chief priests and elders in the temple:

\begin{quote}
A certain man had two sons: and coming to the first, he said: Son, go work to day in my vineyard. And he answering, said: I will not. But afterwards, being moved with repentance, he went. And coming to the other, he said in like manner. And he answering said: I go, Sir. And he went not. Which of the two did the father's will? They say to him: The first. (Matt 21:28--31)
\end{quote}

Neither son begins in serene harmony of feeling and deed. The second speaks with effortless courtesy---\emph{Eo, domine}, ``I go, Sir''---yet leaves the vineyard untouched. The first begins in open refusal, \emph{Nolo} (``I will not''), yet turns back in repentance and puts his hand to the work. St.\ John Chrysostom, expounding this passage in his sixty-seventh \emph{Homily on Matthew}, reads the two sons first as a type of the Gentiles (who had neither promised obedience nor heard the law, yet obeyed in works) and the Jews (who promised at Sinai, ``All that the Lord shall speak, we will do, and will hearken,'' yet disobeyed in works), and then expounds Christ's direct application to the publicans and harlots who repented at John's preaching and obeyed in act.\endnote{St.\ John Chrysostom, \emph{Homilies on the Gospel of St.\ Matthew}, Hom.\ 67.2--3 (NPNF I.10).} Chrysostom emphasizes that the harlots and tax-collectors did not enter the kingdom while remaining what they were, nor by bare profession of sentiment, ``but having obeyed and believed, and having been purified and converted, so did they enter in.''

The parable exposes a quiet falsehood in the claim that reluctant obedience is hypocrisy. Hypocrisy is the second son: polished speech masking refusal of the deed. The first son is initially divided, even defiant, yet his repentance passes into bodily action---he walks into the vineyard. That real act is already a true good, even before the labor feels sweet.

Catholic moral theology safeguards this insight by keeping three distinct levels of human action clear rather than collapsing them into one verdict. St.\ Thomas Aquinas works out the distinction cleanly in his treatise on the Old Law when he asks whether the mode of charity falls under every individual commandment.\endnote{St.\ Thomas Aquinas, \emph{Summa Theologiae} I-II, q.\ 100, a.\ 10, corpus and ad 2--3; cf.\ I-II, q.\ 109, a.\ 4 and q.\ 114, a.\ 2.}

The first level is \emph{objective conformity} of an act to its due norm. Aquinas explains that the commandment ``Honor thy father and thy mother'' requires the substance of the act---honoring one's parents---without making the full interior perfection of charity an automatic condition for avoiding a breach of that specific precept. A person who keeps the fourth commandment, or kneels in Sunday worship, or speaks the truth while his affection is still weak or even before he possesses sanctifying grace does not commit a sin by performing that due act. On the contrary, the act is good in its kind (\emph{bonum ex genere}): it corresponds to what reality and right reason demand. To tell a struggling person that his prayer or worship is worthless unless he feels fervent affection is to contradict the Gospel and to confuse a beginning with a sham.

The second level is \emph{acquired moral virtue}. As Aristotle observed and Aquinas repeats (\emph{ST} I-II, q.\ 107, a.\ 4), it is one thing to do what a just person does, and another to do it as a just person does---steadily, promptly, and with settled ease. Repeatedly performing due acts trains the appetites and the body over time. Duty carried out faithfully while feeling lags behind is the school in which affection itself is educated.

The third level is \emph{supernatural merit in sanctifying grace}. Here Aquinas draws an unyielding line (\emph{ST} I-II, q.\ 100, a.\ 10, ad 2--3; q.\ 114, a.\ 2). Objective conformity to a single commandment and even stable natural virtue cannot by themselves fulfill the supreme commandment to love God with one's whole heart, nor can any human act performed outside sanctifying grace merit eternal life. Eternal beatitude surpasses every created capacity; only an act elevated by divine grace and animated by charity reaches that proportion. Yet this truth never turns imperfect or beginning obedience into sin. Rather, it places human effort in honest relief: the conformed act done before affection matures is genuinely good, genuinely ordered toward God, and genuinely ready to be healed and elevated by grace, while never pretending to be the finished perfection of love.

\subsection{From Body to Community}

St.\ Paul's appeal in Romans 12 confirms why the due return cannot stay locked inside private sentiment. Having asked for bodies presented as a living sacrifice and minds renewed against conformity to the world, the Apostle moves at once to sobriety of self-estimate and membership in one body:

\begin{quote}
For as in one body we have many members, but all the members have not the same office: So we, being many, are one body in Christ; and every one members one of another. (Rom 12:4--5)
\end{quote}

Even at the threshold, the creature's return takes three visible shapes. It is owed to God who owes us nothing, so it cannot work like a commercial exchange. It engages the body and the will in concrete deeds, so it does not wait for mood to authorize obedience. And it belongs to a people joined in one body in Christ, so no individual can discharge it in sealed isolation from his neighbors.
